\section{Proof Details and Technical Notes}
\label{app:proofs}

\subsection{Proof of Theorem 1 (full derivation)}
\label{app:proof_strategy}

The argument is a discrete Lyapunov stability analysis on the coupled
$(P, \Wm)$ system, in the spirit of dwell-time conditions for
switched linear systems~\citep{liberzon2003switching}.  We use
$\Lyap(\Wm_t) = \sum_k w_k \epsilon_k(t)^2$
(Definition~\ref{def:wm_error}) throughout.

\subsubsection*{Verification of C1 and C2$'$ as comparison bounds}

Assumption~C1 decomposes the aggregate correction rate into an
exogenous channel term and an endogenous $\volL$-gradient term.
\emph{Exogenous.}  For each dimension $k$, each channel covering $k$
delivers an EWMA-style correction linear in the residual and
proportional to its sensitivity; summing within $k$ gives a rate
proportional to $\rho_k \cdot \epsilon_k^2$, and bounding below by
$\rhomincross$ across dimensions yields the first term.  Under the
joint-factorization assumption of~\citet{lasser2026exo},
cross-substrate channels' correction signals are conditionally
independent and additive.  \emph{Endogenous.}  The term
$c_{\mathrm{V}} \cdot \nu \cdot \Lyap$ arises from the actor's
$\volL$-gradient: better calibration yields better decisions, which
preserve more cooperative capabilities and raise $\volL$; $\nu$
measures this chain's marginal return along the decision chain
(Assumption~\ref{ass:C1}).

Assumption~C2$'$ supplies the matching upper bound, decomposing
degradation into a Lipschitz term (the world model does not jump
discontinuously on channel loss, giving $c_{\mathrm{D}} \cdot \Lyap$)
and an exogenous drift kick (channels removed from well-calibrated
dimensions drift at a bounded rate $\dmax$, giving
$\dmax \cdot \Delta\rho$).  Proposition~\ref{prop:error_dynamics}
converts C1 and C2$'$ into the per-step affine bound used throughout
the proof.

\subsubsection*{Proof of Part~(a): self-correcting basin}

\paragraph{The affine recursion.}
Proposition~\ref{prop:error_dynamics} gives the per-step bound
\begin{equation*}
\Lyap_{t+1} \leq a_t \cdot \Lyap_t + b_t
\end{equation*}
with
\begin{align*}
a_t &= 1 - \rS(t) \cdot \rhomincross(P_t)
   - \alpha(t) \cdot c_{\mathrm{V}} \cdot \nu(P_t, \Wm_t)
   + \nsub(t) \cdot \rW(t), \\
b_t &= \nsub(t) \cdot \dmax \cdot \Delta\rho(t).
\end{align*}
Iterating over a window $[t_0, t_0 + T]$:
\begin{equation}
\label{eq:iterated_bound}
\Lyap_{t_0 + T} \leq \Lyap_{t_0} \prod_{t=t_0}^{t_0+T-1} a_t
  + \sum_{s=t_0}^{t_0+T-1} b_s \prod_{t=s+1}^{t_0+T-1} a_t.
\end{equation}

\paragraph{The uniform contraction condition.}
For~\eqref{eq:iterated_bound} to yield a bounded limit, we need
aggregate multiplicative decay and bounded additive drift over
sufficiently long windows.  The condition, used in
Theorem~\ref{thm:phase_boundary}(a), is:
\begin{equation}
\label{eq:uniform_contraction}
\exists\, \bar{a} \in (0,1),\ \bar{b} \geq 0,\ T^* > 0:\ \forall T \geq T^*,\quad
\frac{1}{T}\!\!\sum_{t=t_0}^{t_0+T-1}\!\! \ln a_t \leq \ln \bar{a},
\quad
b_t \leq \bar{b} \text{ pointwise}.
\end{equation}
The log-average form is the natural object for $a_t$ because the first
term of \eqref{eq:iterated_bound} is a product; equivalently, the
geometric mean of $a_t$ is at most $\bar{a}$.  The pointwise bound on
$b_t$ is stronger than a time-average bound but is what the
suffix-weighted sum in \eqref{eq:iterated_bound} requires: individual
large $b_s$ terms are not averaged away by the tail product.  A pointwise
bound follows from C2$'$ with $\bar{b} = \dmax \cdot \overline{\Delta\rho}$
(applied at each subsumption step; non-subsumption steps have
$b_t = 0$).  We require $a_t > 0$ throughout, which holds whenever
$\rW(t) \cdot \nsub(t) < 1$ (a step cannot more than wipe out the
current error); this is a weak condition on per-step degradation and
holds under C2$'$ for any finite learning rate.

\paragraph{Convergence under~\eqref{eq:uniform_contraction}.}
Applying AM--GM to the product in~\eqref{eq:iterated_bound},
$\prod_{t=t_0}^{t_0+T-1} a_t \leq \bar{a}^T$, so the first term decays
geometrically.  For the second term, group the tail products by
suffix length: a step with additive kick $b_s$ at position $s$
contributes at most $b_s \cdot \bar{a}^{T-1-s}$ to $\Lyap_{t_0+T}$ in
the worst case (when all subsequent $a_t$ attain the upper bound
$\bar{a}$).  However, individual $a_t$ may exceed $\bar{a}$ on
subsumption steps even though the geometric mean is bounded.  Define
the \emph{overshoot constant}
\begin{equation*}
\mu = \sup_{N \geq 1} \sup_{t_0}\ \frac{1}{\bar{a}^N} \cdot
  \max_{0 \leq s \leq N-1} \prod_{t=t_0+s+1}^{t_0+N} a_t,
\end{equation*}
which is finite whenever an average-dwell-time condition holds: the
number of consecutive subsumption steps (during which $a_t > \bar{a}$)
within any window of length $N$ is at most $\lambda \cdot N + N_0$ for
constants $\lambda \in [0,1)$ and $N_0 \geq 0$.  This is the dwell-time
formulation of~\citet[Ch.~2]{liberzon2003switching} and is stronger
than $\rsub < 1$ (which bounds only asymptotic density).  Under
dwell-time, burst length is uniformly bounded and $\mu < \infty$.
Summing the geometric series:
\begin{equation}
\label{eq:steady_state_bound}
\limsup_{T \to \infty} \Lyap_{t_0+T}
  \leq \frac{\mu \cdot \bar{b}}{1 - \bar{a}}.
\end{equation}

\paragraph{Linearized regime.}
When $|a_t - 1| \ll 1$ for all $t$ (the moderate-step regime of
Remark~\ref{rem:linearized}; cf.\ Equation~\ref{eq:linearized}),
$\ln a_t \approx a_t - 1$, so \eqref{eq:uniform_contraction} is
equivalent to $\frac{1}{T}\sum (1 - a_t) \geq 1 - \bar{a}$.
Substituting the worst-case values $\alphamin = \inf_t \alpha(t)$,
$\underline{\nu} = \inf_t \nu(P_t, \Wm_t)$, and the time-averaged
$\overline{\Delta\rho}$ yields $1 - \bar{a} = \rS \cdot \rhomincross
+ \alphamin \cdot c_{\mathrm{V}} \cdot \underline{\nu} - \rW \cdot
\rsub$.  For the additive term, only subsumption steps contribute
nonzero $b_s$; in the suffix-weighted sum
$\sum_s b_s \prod_{t>s} a_t$, the density of contributing positions
is $\rsub$ and each is bounded pointwise by
$\dmax \cdot \overline{\Delta\rho}$.  The effective $\bar{b}$ in the
linearized steady-state formula is therefore
$\rsub \cdot \dmax \cdot \overline{\Delta\rho}$.  In the linearized
regime, $\mu = 1 + O(\max_t |a_t - \bar{a}|^2)$, which we drop at
the order of the linearization, so~\eqref{eq:steady_state_bound}
reduces to Equation~\ref{eq:neighborhood}.

\paragraph{Policy correctness: $\epsilon_{\mathrm{safe}} > 0$.}
Define
\begin{equation*}
\epsilon_{\mathrm{safe}} = \sup \{\epsilon > 0 :
  \forall \Wm \text{ with } \Lyap(\Wm) \leq \epsilon,\
  \pi^*(\Wm) = \pi^*(\Wm^*)\},
\end{equation*}
where $\pi^*(\Wm)$ is the locally rational policy under world model
$\Wm$.  For a finite action set at each decision point, $\pi^*$ is
determined by the signs of finitely many action gaps $\{G_j\}$.  At
$\Wm^*$, each preserve-optimal agent has $G_j(\Wm^*) > 0$ (the strict
gap assumption); each subsume-optimal agent has $G_j(\Wm^*) < 0$.
Continuity of $G_j$ in $\Wm$ then gives an open neighborhood of
$\Wm^*$ on which every sign is preserved, hence
$\epsilon_{\mathrm{safe}} > 0$.

A computable lower bound follows from Cauchy--Schwarz in the
$w$-weighted metric.  Writing
$G_j(\Wm) - G_j(\Wm^*) = \int_0^1 \nabla G_j(\Wm_\tau) \cdot
(\Wm - \Wm^*) \, d\tau$
along any path $\Wm_\tau$ from $\Wm^*$ to $\Wm$,
Cauchy--Schwarz gives
$|G_j(\Wm) - G_j(\Wm^*)| \leq \|\nabla G_j\|_{w^{-1}} \cdot \sqrt{\Lyap(\Wm)}$
where $\|\cdot\|_{w^{-1}}$ is the dual norm with weights $1/w_k$.
An action flip on agent $j$ requires
$|G_j(\Wm) - G_j(\Wm^*)| \geq G_j(\Wm^*)$, so:
\begin{equation}
\label{eq:eps_safe_bound}
\epsilon_{\mathrm{safe}} \geq
  \min_{j \in \mathcal{A}^+}
  \left( \frac{G_j(\Wm^*)}
              {\sup_{\Wm \in \mathcal{N}} \|\nabla G_j(\Wm)\|_{w^{-1}}} \right)^2,
\end{equation}
where $\mathcal{A}^+$ is the set of preserve-optimal agents and
$\mathcal{N} = \{\Wm : \Lyap(\Wm) \leq \Lyap_\infty\}$ is the
steady-state neighborhood from~\eqref{eq:steady_state_bound}.  The
self-consistency is that $\Lyap_\infty$ enters both the bound and
$\mathcal{N}$; in practice one takes the smallest $\mathcal{N}$
containing the trajectory and iterates the self-consistency condition
$\Lyap_\infty < \epsilon_{\mathrm{safe}}(\mathcal{N})$ to closure.
The bound~\eqref{eq:eps_safe_bound} is conservative because
$\|\nabla G_j\|_{w^{-1}}$ is taken as its supremum over the entire
neighborhood, not the value along the actual trajectory.  A cheaper
per-coordinate estimator $\min_k w_k \cdot \epsilon_{\mathrm{crit},k}^2$
rules out single-coordinate flips but does not bound multi-coordinate
flips (where the action gap responds to error on several dimensions
simultaneously), so~\eqref{eq:eps_safe_bound} is the correct object
for policy correctness on the whole neighborhood.

\paragraph{Closing Part~(a).}
Combining~\eqref{eq:steady_state_bound} with the policy-correctness
condition $\Lyap_\infty < \epsilon_{\mathrm{safe}}$:  once the
trajectory enters the $\epsilon_{\mathrm{safe}}$-neighborhood of
$\Wm^*$, the locally rational policy agrees with the true
$\Vdisc$-optimal policy and the anti-monopolar
result~\citep[Proposition~6]{lasser2026horizon} applies.  As
subsumption ceases ($\overline{\Delta\rho} \to 0$), $\bar{b} \to 0$
and the neighborhood shrinks to~$\{0\}$.

\paragraph{Co-evolution caveats.}
Two dependencies complicate the fixed-parameter reading.  First, the
contraction factor $a_t$ depends on $\rhomincross(P_t)$, which evolves
as subsumption proceeds; the self-correcting basin is defined over
trajectories, not initial states.  Second, $\alpha(t)$ depends on
$T_s$~\citep{lasser2026gfm}, which depends on the error history; the
EWMA structure of $T_s$ makes $\alpha$ monotone in prediction
residuals, so $\rS$ is bounded below by a function of the minimum
$T_s$ along the trajectory.  The uniform contraction
condition~\eqref{eq:uniform_contraction} is the object that survives
both dependencies because it is stated in terms of time averages of
the actual $a_t$, not instantaneous parameter values.

\subsubsection*{Proof of Part~(b): absorbing basin}

The triggering conditions are (i) $B(t_0, T) > 0$
(Equation~\ref{eq:error_budget}), (ii) $\rhocross_k = 0$ on a cascade
dimension $k$, and (iii) $\epsilon_k > \epsilon_{\mathrm{crit},k}$.
Under B1a, (ii) gives $\epsilon_k$ non-decreasing; environmental
non-stationarity on $k$ then keeps (iii) satisfied.  The Cascade
Lemma~\ref{lem:cascade} converts miscalibration on $k$ into action
flips on agents covering dimensions $j \neq k$; each resulting
subsumption drives $\rhocross_j \to 0$ on at least one new dimension
(agent-channel graph connectivity) and the per-coordinate cascade
threshold $\epsilon_{\mathrm{crit},j}$ is crossed there too under
non-stationarity.  Iterating across dimensions, the trajectory
converges to a neighborhood of $S^*$
(Definition~\ref{def:monopolar}).  The per-coordinate cascade
threshold is the correct object here because Part~(b) concerns
propagation through dimensions sequentially; the aggregate
$\epsilon_{\mathrm{safe}}$ bound of~\eqref{eq:eps_safe_bound} is
a sufficient condition for the entire neighborhood to be safe, while
$\epsilon_{\mathrm{crit},k}$ is the necessary threshold for the
\emph{specific} action flip on dimension~$k$.

\subsubsection*{Proof of Part~(b$'$): metastable basin}

Under B1$'$ (Assumption~\ref{ass:B1prime}), the indirect signal $g_k$
provides residual correction at information rate $I_k$, so
$\epsilon_k$ is no longer non-decreasing once the trajectory enters
the absorbing basin of Part~(b).  The cascade-trigger analysis of
Part~(b) still applies (B1$'$ does not remove the triggering
mechanism, only the irreversibility), but escape from the basin is
now possible through $g_k$.  The lifetime $\tau_{\mathrm{meta}}$ is
set by the competition between the basin depth and the information
rate.  Proposition~\ref{prop:metastable} gives the scaling law
$\tau_{\mathrm{meta}} \gtrsim C/I_k$, where the constant $C$ depends
on the basin geometry near $S^*$ and is not characterized as a
theorem-level bound.  The proof of Proposition~\ref{prop:metastable}
itself (by the mean first-passage argument of Kramers-type escape)
is deferred.

\subsubsection*{Proof of Part~(c): critical surface}

Setting $B(t_0, T) = 0$ in the linearized regime of
Part~(a)---equivalently, requiring $(1 - \bar{a}) = 0$---gives
$\rS \cdot \rhomincross + \alphamin \cdot c_{\mathrm{V}} \cdot
\underline{\nu} = \rW \cdot \rsub + \dmax \cdot
\Delta\rho_{\mathrm{avg}} / \Lyap_{\mathrm{avg}}$.  Solving for
$\rhomincross$ under the standing assumption $\rS > 0$ yields
Equation~\ref{eq:critical_redundancy}.  The surface is linearized: the
nonlinear generalization, where $\rS$, $\rW$, and $\rsub$ depend on
$(P_t, \Wm_t)$, would produce a state-dependent critical surface
rather than a hyperplane (see Section~\ref{sec:discussion}).

\subsection{Cascade Lemma}
\label{app:cascade}

\begin{lemma}[Cascade]
\label{lem:cascade}
Define the \emph{action gap} for agent $a_j$ at world-model state
$\Wm$ as
\begin{equation*}
G_j(\Wm) = \Vdisc^{\mathrm{preserve}}(a_j, \Wm)
  - \Vdisc^{\mathrm{subsume}}(a_j, \Wm)
\end{equation*}
(positive when preservation is preferred).  Fix a dimension $k$ and consider a state $\Wm_t$ satisfying:
\begin{enumerate}
\item[(0)] \textbf{Single-coordinate dominance.}
  $\Wm_t[k'] = \Wm^*[k']$ for all $k' \neq k$.  (Dimension $k$ has
  been blinded while all other dimensions are still corrected.  This
  is the cascade scenario: the lemma is applied at the step where $k$
  first triggers a wrong action, before the resulting subsumption
  blinds further dimensions.)
\end{enumerate}

\noindent Define the one-dimensional path
$\Wm(\epsilon) = \Wm^* + \epsilon \cdot e_k \cdot
\mathrm{sgn}(\Wm_t[k] - \Wm^*[k])$, where $e_k$ is the $k$-th
standard basis vector and $\epsilon \in [0, \epsilon_k(t)]$.  Under
assumption~(0), $\Wm_t = \Wm(\epsilon_k(t))$.

Assume additionally:
\begin{enumerate}
\item[(i)] $G_j(\Wm(\epsilon))$ is continuously differentiable in
  $\epsilon$,
\item[(ii)] the miscalibration direction on $k$ makes preservation
  less attractive:
  $\frac{d}{d\epsilon} G_j(\Wm(\epsilon)) < 0$ for all
  $\epsilon \in [0, \epsilon_k(t)]$ (i.e., increasing error on $k$
  monotonically decreases the action gap),
\item[(ii$'$)] \textbf{(Multi-dimensional monotonicity; needed for
  cascade propagation beyond the first trigger.)}  For every
  safety-relevant dimension $k'$ and every agent $a_{j'}$,
  $\sign(\Wm[k'] - \Wm^*[k']) \cdot
  \frac{\partial G_{j'}}{\partial \Wm[k']} < 0$ whenever
  $\Wm[k'] \neq \Wm^*[k']$; that is, moving $\Wm[k']$ further
  from $\Wm^*[k']$ (increasing the absolute error $\epsilon_{k'}$)
  monotonically decreases every agent's preservation advantage,
  regardless of the sign of the miscalibration.
  Assumption~(ii) is the single-coordinate
  special case (where the error direction is fixed by construction
  of the path $\Wm(\epsilon)$); (ii$'$) extends it to the
  multi-coordinate regime where multiple dimensions are simultaneously
  miscalibrated with potentially different error signs.
\end{enumerate}

Define the cascade threshold for agent $a_j$ on dimension $k$:
\begin{equation}
\label{eq:cascade_threshold_def}
\epsilon_{\mathrm{crit}}^{(j,k)} =
  \frac{G_j(\Wm^*)}
  {\displaystyle\inf_{\epsilon \in [0, \epsilon_k(t)]}
    \left|\frac{d}{d\epsilon} G_j(\Wm(\epsilon))\right|}
\end{equation}
(the true action gap divided by the minimum sensitivity to
dimension-$k$ error along the path; a lower bound on sensitivity
gives an upper bound on the error needed to flip the action).  The
per-dimension cascade threshold is $\epsilon_{\mathrm{crit},k} =
\min_{j \in \mathcal{A}_t^+} \epsilon_{\mathrm{crit}}^{(j,k)}$,
minimized over the set of \emph{preserve-optimal agents}
$\mathcal{A}_t^+ = \{j \in \mathcal{A}_t : G_j(\Wm^*) > 0\}$.
Agents that are truly subsume-optimal ($G_j(\Wm^*) \leq 0$) are
excluded because their subsumption is correct regardless of
$\epsilon_k$; only flips from correct-preserve to
incorrect-subsume constitute cascade triggers.

Then the cascade operates in two steps:
\begin{enumerate}
\item[\textbf{(A)}] \textbf{Action flip.}  If
  $\epsilon_k(t) > \epsilon_{\mathrm{crit}}^{(j,k)}$, then $G_j(\Wm_t) < 0$
  and the locally rational action under $\Wm_t$ selects subsumption of
  $a_j$.
\item[\textbf{(B)}] \textbf{Error propagation.}  Subsumption of $a_j$
  removes $a_j$'s channels.  For any dimension $j' \neq k$ covered
  exclusively by $a_j$'s channels, the per-coordinate error
  $\epsilon_{j'}(t+1) \geq \epsilon_{j'}(t)$.
\end{enumerate}

Miscalibration on one dimension produces decisions that degrade other
dimensions.
\end{lemma}

\begin{proof}
\emph{Step (A): Action flip.}
Under the true world model $\Wm^*$, preservation is optimal:
$G_j(\Wm^*) > 0$.  By the mean value theorem applied to the
scalar function $\epsilon \mapsto G_j(\Wm(\epsilon))$ on
$[0, \epsilon_k(t)]$:
\begin{equation*}
G_j(\Wm_t) = G_j(\Wm(\epsilon_k(t)))
  = G_j(\Wm^*) +
  \frac{d G_j}{d\epsilon}\bigg|_{\epsilon = \tilde{\epsilon}}
  \cdot \epsilon_k(t)
\end{equation*}
for some $\tilde{\epsilon} \in (0, \epsilon_k(t))$.  By
assumption~(ii), the derivative is negative, so:
\begin{equation*}
G_j(\Wm_t) = G_j(\Wm^*)
  - \left|\frac{d G_j}{d\epsilon}\bigg|_{\tilde{\epsilon}}\right|
  \cdot \epsilon_k(t)
  \leq G_j(\Wm^*)
  - \inf_{\epsilon} \left|\frac{d G_j}{d\epsilon}\right|
  \cdot \epsilon_k(t)
\end{equation*}
When $\epsilon_k(t) > \epsilon_{\mathrm{crit}}^{(j,k)} =
G_j(\Wm^*) \big/ \inf |\frac{d G_j}{d\epsilon}|$, the right side is
negative, so $G_j(\Wm_t) < 0$ and the locally rational action
selects subsumption of $a_j$.

\emph{Step (B): Error propagation.}
Once $a_j$ is subsumed, Lemma~\ref{lem:channel_contraction} gives
$\Obs(P_{t+1}) \subseteq \Obs(P_t)$.  For any dimension $j'$
covered exclusively by $a_j$'s channels, $\rhocross_{j'}(P_{t+1})
= 0$.  By Assumption~B1a (no endogenous correction on blind dimensions;
Assumption~\ref{ass:B1}), $f_{\Wm}$ does not decrease
$\epsilon_{j'}$ when $\rhocross_{j'} = 0$.  Since the correction
signal vanishes, $\epsilon_{j'}(t+1) \geq \epsilon_{j'}(t)$
(non-decreasing).  Under environmental non-stationarity on
dimension~$j'$ (the true state continues to evolve while the blind
estimate cannot track it), $\epsilon_{j'}$ grows strictly; C2$'$
bounds the growth rate by $\dmax$ but the cascade requires only that
drift is positive, not that it saturates the bound.

The cascade propagates: if $\epsilon_{j'}$ subsequently exceeds
$\epsilon_{\mathrm{crit},j'}$ for another agent $a_{j''}$ whose channels
cover dimensions $j'' \neq j'$, step~(A) applies again, driving
further subsumptions.

\emph{Multi-coordinate reinforcement at later stages.}
After the first trigger, assumption~(0) is violated for later
cascade stages: dimensions $k$ and $j'$ are both miscalibrated.
However, the errors on previously-blinded dimensions all push $G_j$
in the same direction (more negative, making subsumption more
attractive), provided assumption~(ii) holds uniformly across all
dimensions, not just the original trigger dimension~$k$.  This
requires a multi-dimensional monotonicity condition: for every
previously-blinded dimension $k_i$ and every agent $a_{j_n}$ at
cascade stage~$n$,
$\sign(\Wm[k_i] - \Wm^*[k_i]) \cdot
\frac{\partial G_{j_n}}{\partial \Wm[k_i]} < 0$.
This is assumption~(ii$'$) above, which is stated as an explicit
condition of the lemma.  It is plausible for the
subsumption-vs-preservation action gap (increasing error on any
safety-relevant dimension decreases preservation value) and holds
when miscalibration uniformly favors subsumption.  Under~(ii$'$), at
cascade stage $n$ with errors $\epsilon_{k_1}, \ldots,
\epsilon_{k_{n-1}}$ on previously blinded dimensions, the action gap
satisfies
\begin{equation*}
G_{j_n}(\Wm_t) \leq G_{j_n}(\Wm^*) - \sum_{i=1}^{n}
  \inf \left|\frac{\partial G_{j_n}}{\partial \Wm[k_i]}\right|
  \cdot \epsilon_{k_i}
\end{equation*}
(by iterated application of the MVT along each coordinate
successively).  Since every term in the sum is non-negative, the
multi-coordinate case is \emph{easier} to trigger than the single-
coordinate case: the effective threshold for the $n$-th trigger
is lower, not higher.  The single-coordinate threshold
$\epsilon_{\mathrm{crit}}^{(j_n,k_n)}$ is therefore a conservative upper bound
at every cascade stage.
\end{proof}

\subsection{Connection to the discount-factor threshold $\gamma^*$}
\label{app:gamma_star}

The anti-monopolar result~\citep[Proposition~6]{lasser2026horizon} gives $\gamma^* = 1 - \rext / \Delta_0$ as the
critical discount factor for diversity dominance.  Under the coupled
$(P, \Wm)$ dynamics, the actor does not know the true $\rext$; it uses
$\hat{\rext} = \rext - \delta_{\rext}$.  The \emph{perceived}
critical discount factor is:
\begin{equation}
\label{eq:perceived_gamma}
\gamma^*_{\mathrm{perceived}} = 1 -
  \frac{\rext - \delta_{\rext}}{\Delta_0}
\end{equation}

When $\delta_{\rext} > 0$ (underestimating cross-substrate growth):
\begin{itemize}
\item $\gamma^*_{\mathrm{perceived}} > \gamma^*$ (the threshold
  appears higher than it truly is),
\item For $\gamma \in (\gamma^*, \gamma^*_{\mathrm{perceived}})$,
  diversity dominates in truth but domination appears rational to the
  actor.
\end{itemize}

This is the \emph{mechanism} by which world-model error converts a
self-balancing regime into a locally rational subsumption cascade.  The
phase boundary from Theorem~\ref{thm:phase_boundary} characterizes
when $\delta_{\rext}$ stays small enough that
$\gamma^*_{\mathrm{perceived}}$ remains close to $\gamma^*$,
preserving the effective reach of the anti-monopolar result~\citep{lasser2026horizon}.
